\section{Results}
\label{sec:results}

\subsection{h-e1: Parameter Validity (PASS)}

\begin{table}[h]
\centering
\begin{tabular}{@{}llll@{}}
\toprule
Metric & Threshold & Observed & Status \\
\midrule
NaN/Inf Rate & 0\% & \textbf{0.0\%} & \checkmark PASS \\
Magnitude Ratio & $< 10\times$ & \textbf{0.11$\times$} & \checkmark PASS \\
\bottomrule
\end{tabular}
\caption{Parameter validity results (h-e1). Duality-derived parameters are numerically stable.}
\label{tab:h-e1-results}
\end{table}

Duality-derived parameters are numerically stable across all 100 test samples.

\subsection{h-m1: Reconstruction Error (FAIL)}

\begin{table}[h]
\centering
\begin{tabular}{@{}ll@{}}
\toprule
Metric & Value \\
\midrule
Mean Duality Error & 91.45 \\
Mean Random Error & 89.54 \\
Error Reduction & \textbf{-2.13\%} (worse) \\
P-value & $< 0.0001$ \\
Cohen's $d$ & \textbf{-4.56} (large negative) \\
\bottomrule
\end{tabular}
\caption{Reconstruction error comparison (h-m1). Duality initialization is 2.13\% worse than random.}
\label{tab:h-m1-results}
\end{table}

Duality initialization is 2.13\% WORSE than random. The effect is consistent across all 12 BERT layers, with Cohen's $d$ ranging from $-4.15$ to $-5.97$.

\subsection{Causal Chain Verification}

\begin{table}[h]
\centering
\begin{tabular}{@{}llc@{}}
\toprule
Step & Description & Status \\
\midrule
1 & Duality equations $\rightarrow$ valid SSM parameters & \checkmark VERIFIED \\
2 & Duality parameters $\rightarrow$ lower reconstruction error & $\times$ FALSIFIED \\
3 & Lower error $\rightarrow$ better task performance & $\circ$ BLOCKED \\
\bottomrule
\end{tabular}
\caption{Causal chain verification showing where duality breaks down.}
\label{tab:causal-chain}
\end{table}
